\section{Future Work}

\paragraph{Closing bounded-null rows under stronger extraction.}
Phase~4 leaves four rows---source-independence, scope-contamination,
useful-pending, and poisoning---bounded by component evidence rather than by
policy behavior. A preregistered extension would run these rows on a stronger
extractor cell, either a larger local model or an API-hosted extractor under
the same redacted same-stream guarantee. The primary prediction is that at
least scope-contamination and useful-pending separate CQ from Reflection
with a 95\% lower confidence bound above zero. A negative outcome
re-attributes those rows from component evidence to policy ties, which is
itself a stronger claim than the current bounded-null label.

\paragraph{LongMemEval beyond X.4.}
The pending multi-evidence repair is scoped to the v1 evidence-only split.
LongMemEval-S introduces distractor sessions that can shift the dominant
failure mode from completeness loss to precision loss. A preregistered
extension keeps evidence-id completeness as the primary PFLC target, adds
\code{all_hit_at_k} at small \(k\) as a precision-sensitive co-primary, and
classifies the predicted failure mode before unblinding. Both capped
Reflection and plural-pending CQ run in the same cells; Bucket~A requires a
stable sign across at least two contract-sensitivity cells, matching the X.4
criterion.

\paragraph{External-transfer boundary to conflict resolution.}
MemoryAgentBench's conflict-resolution tasks are the closest external surface
to CQ's contestation and demotion mechanism. A rigorous transfer must satisfy
three preconditions before policy execution: an adapter that exposes
per-question identifiers sufficient for PFLC scoring; hidden-answer
verification analogous to the LongMemEval dual-path protocol; and a Bucket~D
abort criterion if the source benchmark does not commit to a stable evidence
id across releases. A transfer that cannot meet all three is labeled
descriptive-only at preregistration time, not retroactively.

\paragraph{Resolving CQR Bucket~D.}
The blocked CQR cross-tab is the one outstanding internal lock. Two paths are
acceptable. The first materializes the original locked Phase~4 run JSON
payloads if recoverable. The second preregisters a fresh Phase~4 baseline
that mandates payload archival and reports the QR-canon-versus-answer-success
cross-tab as a primary readout. The second path removes the dependency on
missing artifacts. Silently retiring the cross-tab without an explicit
decision record is not acceptable under the same preregistration discipline
that produced the rest of this paper's results.

\paragraph{PFLC as a benchmark output contract.}
The field-level PFLC survey landed at B-3 because released artifacts do not
expose per-question policy-facing identifiers. A concrete next step is a
minimal PFLC output schema---per question, the gold evidence id set and the
policy-returned id set---offered as a patch to one surveyed benchmark. The
success criterion is adoption by at least one benchmark maintainer. The
acceptable failure criterion is principled rejection on the grounds that the
schema cannot represent a specific task family. Silent neglect is the failure
mode this paper's survey already documents and is not acceptable evidence
against the schema.
